\documentclass[11pt]{article}

\usepackage[T1]{fontenc}
\usepackage[utf8]{inputenc}
\usepackage{lmodern}
\usepackage{microtype}
\usepackage[a4paper,margin=27mm,headheight=15pt]{geometry}
\usepackage{amsmath,amssymb,mathtools}
\usepackage{booktabs,tabularx,longtable,array}
\usepackage{enumitem}
\usepackage{xcolor}
\usepackage{hyperref}
\usepackage{fancyhdr}
\usepackage{titlesec}

\definecolor{navy}{HTML}{17365D}
\definecolor{bluegray}{HTML}{506784}
\definecolor{lightblue}{HTML}{EAF1F8}
\definecolor{lightgreen}{HTML}{EAF5EE}
\definecolor{lightred}{HTML}{FBEDEE}
\definecolor{lightamber}{HTML}{FFF5DC}
\definecolor{darkgreen}{HTML}{21633B}
\definecolor{darkred}{HTML}{8A2633}
\definecolor{darkamber}{HTML}{77550B}

\hypersetup{
  colorlinks=true,
  linkcolor=navy,
  urlcolor=bluegray,
  pdftitle={Cold Attack, Round 2: Borrowed-Edge Countermodels and the Anchor Residue},
  pdfsubject={Round-2 attack on the mixed quadrant of the forall-PO-free ALCI-RCC5 fragment}
}

\pagestyle{fancy}
\fancyhf{}
\fancyhead[L]{\small\color{bluegray}Mixed quadrant, round 2}
\fancyhead[R]{\small\color{bluegray}Independent technical attack}
\fancyfoot[C]{\small\thepage}
\renewcommand{\headrulewidth}{0.4pt}

\titleformat{\section}{\Large\bfseries\color{navy}}{\thesection}{0.65em}{}
\titleformat{\subsection}{\large\bfseries\color{navy}}{\thesubsection}{0.65em}{}
\titleformat{\subsubsection}{\normalsize\bfseries\color{bluegray}}{\thesubsubsection}{0.65em}{}
\setlist[itemize]{leftmargin=1.6em,itemsep=0.25em,topsep=0.35em}
\setlist[enumerate]{leftmargin=1.8em,itemsep=0.25em,topsep=0.35em}
\setlength{\parindent}{0pt}
\setlength{\parskip}{0.55em}
\allowdisplaybreaks

\newcommand{\PP}{\mathsf{PP}}
\newcommand{\PPI}{\mathsf{PPI}}
\newcommand{\PO}{\mathsf{PO}}
\newcommand{\DR}{\mathsf{DR}}
\newcommand{\EQ}{\mathsf{EQ}}
\newcommand{\mty}{\operatorname{mty}}
\newcommand{\comp}{\operatorname{comp}}
\newcommand{\cl}{\operatorname{cl}}
\newcommand{\dkey}{\operatorname{dkey}}
\newcommand{\statusbox}[2]{%
  \begin{center}
  \fcolorbox{#1}{#1!8}{\begin{minipage}{0.91\linewidth}\textbf{#2}\end{minipage}}
  \end{center}}
\newcommand{\evP}{\textbf{[P]}}
\newcommand{\evX}{\textbf{[X]}}
\newcommand{\evS}{\textbf{[S]}}
\newcommand{\evR}{\textbf{[R]}}
\newcommand{\evO}{\textbf{[O]}}

\begin{document}

\pagenumbering{gobble}
\begin{titlepage}
\thispagestyle{empty}
\vspace*{1.35cm}
{\color{navy}\rule{\linewidth}{1.2pt}}\\[1.0cm]
{\huge\bfseries\color{navy}Cold Attack, Round 2\\[0.24cm]
Borrowed-Edge Countermodels\\[0.12cm]
and the Anchor Residue\par}
\vspace{0.65cm}
{\Large\color{bluegray}The mixed quadrant of the $\forall PO$-free fragment of
$\mathcal{ALCI}_{\mathrm{RCC5}}$\par}
\vspace{1.15cm}

\begin{center}
\begin{tabular}{rl}
\textbf{Status:} & Independent technical attack, not a referee report \\
\textbf{Input packet:} & \texttt{attack\_mixed\_quadrant\_r2.zip} \\
\textbf{Report date:} & 2026-08-28 \\
\textbf{Evidence:} & Exact finite models, symbolic arguments, and executable checks
\end{tabular}
\end{center}

\vfill
\noindent\fcolorbox{darkred}{lightred}{%
\begin{minipage}{0.94\dimexpr\linewidth-2\fboxsep-2\fboxrule\relax}
\textbf{Principal result.}
Both open halves of Target D are false for the current borrowed-edge discipline.
There is a refined-key blocked node whose gate-mate has no witness avoiding
disjointness (D2), and there is a separate exact-key construction in which every
fallback pair is $\PO$ but two groups form a directed cycle (D1). Target E splits:
finite keyed node generation terminates, but the current gate-mate borrowing rule is
not anchor-compatible and need not extend to a valid kernel frame.
\end{minipage}}
\par
\vspace{0.85cm}
{\color{navy}\rule{\linewidth}{1.2pt}}
\end{titlepage}

\pagenumbering{roman}
\begin{abstract}
The round-2 packet reports substantial progress since the first attack: finite
cross-kernel fusion is now Lean-formalized, and the down-spectrum gate proposed in
round 1 has been integrated. The remaining construction borrows a gate-mate's
witness for a blocked node. This report refutes both safety claims needed by that
discipline. First, a seven-point mixed set model gives an actual refined-key blocked
node for which the gate-mate's unique $\PP$ witness is $\DR$ from the borrower; the
declared order stays acyclic but violates \texttt{ltNotDj}. Second, a 14-point set
model gives two exact-key fallback groups whose borrowed pairs are all $\PO$ and
individually safe, yet model-$\PP$ return paths make the declared order cyclic.

For Target E, the combinatorial node closure does terminate after at most one
addition per key-demand pair. That finite external fixed point is not enough: a
symmetric two-component model shows that gate-mate borrowing can declare two anchors
disjoint although a selected kernel point lies below both, so no compatible kernel
placement exists. The report also audits the supplied probes. In particular,
\texttt{wp131} still uses the old type-only gate, aliases children across roles, and
confounds agreement with repeat-free selection; \texttt{wp132}'s fresh-point
evaluator is not a sound satisfiability test. Candidate repairs are isolated, but no
complete extraction or decidability theorem is claimed.
\end{abstract}

\tableofcontents
\clearpage
\pagenumbering{arabic}

\section{Executive verdict and evidence ledger}

The input packet distinguishes two borrowed-edge obligations and one anchor-closure
question. The cold attack gives the following answers.

\begin{center}
\renewcommand{\arraystretch}{1.3}
\begin{tabularx}{\linewidth}{>{\bfseries}p{0.20\linewidth}X>{\raggedright\arraybackslash}p{0.23\linewidth}}
\toprule
Question & Finding & Status \\
\midrule
D1: global acyclicity & False for the local agreement-first/fallback discipline.
Two individually safe $\PO$ fallback groups can alternate into a cycle. & Exact finite
set-model counterexample \evX; a globally coordinated assignment remains possible. \\
D2: safe existing witness & False. At the exact down-spectrum key, the gate-mate's
only relevant witness can be $\DR$ from the blocked node. & Exact mixed
counterexample \evX; independent exhaustive core search. \\
E: finite closure & Pure keyed node generation terminates with the advertised finite
bound. & Combinatorial lemma \evP/\evX, essentially the packet's abstract gate theorem. \\
E: valid placement & The current gate-mate borrowing policy can produce a finite
external frame that has no compatible kernel extension. & Exact finite counterexample
\evX; general anchor-compatible completion remains open \evO. \\
Probe evidence & Several headline measurements do not test the stated construction,
and one wrapper masks failures. & Deterministic audit \evX; randomized counts remain
measurements only \evR. \\
\bottomrule
\end{tabularx}
\end{center}

\subsection{Evidence levels}

The tags used throughout are deliberately strict.

\begin{itemize}
  \item \evP{} is a theorem reported by the packet as Lean-certified. The archive does
        not contain the Lean source, theorem statements, build manifest, or compiler
        output, so it was not independently rebuilt here.
  \item \evX{} is an exact finite model or finite combinatorial argument accompanied by
        a deterministic executable checker. It is mathematical evidence, but not a Lean
        proof.
  \item \evS{} is a symbolic infinite family, with bounded executable regression checks.
  \item \evR{} is a randomized finite sweep. It can find counterexamples but cannot prove
        a universal statement.
  \item \evO{} is an open construction or candidate reduction.
\end{itemize}

The packet reports a roughly 42,700-line Lean development with no sorries or warnings
and axioms \texttt{propext}, \texttt{Quot.sound}, and
\texttt{Classical.choice}. Those claims are treated as premises, not as independently
reviewed results.

\subsection{Continuity with round 1}

Round 1's finite-rectangle result is reported formalized end to end as
\texttt{atom\_seq\_stabilizes}, \texttt{two\_tower\_rectangle\_gen},
\texttt{finite\_fusion\_recurrent}, and \texttt{fused\_kq\_all}. The packet also
reports that the type-only common-top counterexample was accepted and that the
down-spectrum key and local cone-transfer lemmas were certified. This report does not
re-derive those results. It attacks the new claims that are logically downstream from
them.

\section{Exact setting and the quantifiers at issue}

Let $\rho$ be an atomic RCC5 model on a finite node set $V$. Strong equality is used:
$\rho(x,y)=\EQ$ exactly when $x=y$. The convention is

\[
  \rho(x,y)=\PP \quad\Longleftrightarrow\quad x\text{ is a proper part of }y.
\]

The refined gate key is
\[
  q(v)=\left(\mty(v),\;\{\mty(x):\rho(x,v)=\PP\}\right).
\]
The gate retains the first node of each key. If $v$ is blocked and $a$ is its retained
gate-mate, then $q(a)=q(v)$. For a nonpersistent demand
$\delta=\exists\PP.D\in\mty(v)$, the current construction chooses one gate-mate
witness $z$ with
\[
  \rho(a,z)=\PP,\qquad z\models D,
\]
and declares a borrowed order edge $v<z$.

The declared order seed is
\[
  L=\{(x,y)\in V^2:\rho(x,y)=\PP\}\cup B,
\]
where $B$ is the set of borrowed edges. The frame must be ordered-disjoint:

\begin{itemize}
  \item the transitive closure $L^+$ is irreflexive;
  \item \texttt{disj} is symmetric, irreflexive, and downward closed in both
        coordinates;
  \item $x<y$ implies $\neg\operatorname{disj}(x,y)$.
\end{itemize}

The packet's certified dichotomy concerns one edge: if a chosen $z$ lies below $v$,
then $a$ lies below $v$, where an agreeing witness exists. This yields a local choice
with $\rho(z,v)\neq\PP$. It does not settle either of the following:

\begin{description}
  \item[D1.] Do several locally safe borrowed edges keep $L^+$ irreflexive?
  \item[D2.] In the non-agreeing case, can $z$ always be chosen with
        $\rho(v,z)\neq\DR$ as well?
\end{description}

These are independent. D2 is a pairwise compatibility condition. D1 is a global graph
condition, even if every borrowed pair is $\PO$ and hence already satisfies D2.

\section{Target D2: the safe existing witness need not exist}

\statusbox{darkred}{Result \evX: D2 is false at the exact refined key. The failure
uses a genuinely mixed, $\forall PO$-free input, an actual blocking run, a
nonpersistent demand, a unique gate-mate witness, and equal nonempty lower spectra.}

\subsection{The seven-point mixed model}

Use ordinary nonempty-set RCC5 semantics on the following points. Atomic propositions
not displayed are false.

\begin{center}
\renewcommand{\arraystretch}{1.18}
\begin{tabular}{ccl}
\toprule
Point & Region & True atoms \\
\midrule
$a_0$ & $\{0\}$ & --- \\
$a$   & $\{0,1\}$ & $Q$ \\
$z$   & $\{0,1,2\}$ & $P$ \\
$v_0$ & $\{3\}$ & --- \\
$v$   & $\{3,4\}$ & $Q$ \\
$w$   & $\{3,4,5\}$ & $P$ \\
$p$   & $\{1,4\}$ & $R$ \\
\bottomrule
\end{tabular}
\end{center}

Let
\[
  C_0=(\exists\DR.Q)\sqcap(\exists\PP.P)\sqcap(\exists\PO.R).
\]
It contains no universal $\PO$ formula and lies in the mixed existential quadrant.
Both $a$ and $v$ satisfy $C_0$:

\begin{itemize}
  \item they are each other's unique $\DR.Q$ witnesses;
  \item $z$ is the unique $\PP.P$ witness of $a$, and $w$ is the unique
        $\PP.P$ witness of $v$;
  \item $p$ is a $\PO.R$ witness of both.
\end{itemize}

The component swap
\[
  a_0\leftrightarrow v_0,\qquad a\leftrightarrow v,\qquad z\leftrightarrow w,
  \qquad p\leftrightarrow p
\]
is a valuation-preserving automorphism. Therefore
$\mty(a)=\mty(v)$ over $\cl(C_0)$. Moreover, the strict lower cones of $a$ and $v$
are $\{a_0\}$ and $\{v_0\}$, and $a_0,v_0$ have the same closure type. Hence
\[
  q(a)=q(v),
\]
with equal lower spectra of cardinality one. The example does not rely on an empty
spectrum.

\subsection{The actual gate run}

Start with the duplicate-free singleton seed $[a]$. Its three existential demands have
unique witnesses, so the first extension adds
\[
  v\quad(\DR.Q),\qquad z\quad(\PP.P),\qquad p\quad(\PO.R).
\]
The refined gate scans $[a,v,z,p]$, keeps $a$ first, and blocks $v$ because
$q(a)=q(v)$. The $\PP.P$ demand is not kernel-routed: its persistence guard
\[
  \forall\PP.(\exists\PP.P)
\]
is false at $a$ because the maximal point $z$ has no $\PP.P$ successor.

The complete gate-mate candidate set is the singleton
\[
  \{y:\rho(a,y)=\PP\text{ and }y\models P\}=\{z\}.
\]
But the two components are disjoint, so
\[
  \rho(v,z)=\rho(z,v)=\DR.
\]
Consequently
\[
  \{y:\rho(a,y)=\PP,\ y\models P,\ \rho(y,v)\neq\PP,
        \ \rho(v,y)\neq\DR\}=\varnothing.
\]

Declaring $v<z$ remains acyclic in this small model, which isolates D2 from D1.
However, inherited disjointness contains $v\mathrel{\operatorname{disj}}z$, so the
frame clause $v<z\Rightarrow\neg\operatorname{disj}(v,z)$ fails immediately.

\subsection{What the counterexample rules out}

The witness set is a singleton. Greedy, maximal, group-aware, and any other selection
rule restricted to the gate-mate's existing witnesses make the same choice. Reversing
which twin occurs first fails symmetrically. Thus the counterexample refutes the
universal D2 existence claim, not merely one implementation of selection.

It does \emph{not} refute expanding $v$ and using its own witness $w$, adding a fresh
class gadget, refining the key with additional incidence information, or changing the
global blocking discipline. Those alternatives need new bounds and certificate proofs.
The packaged exhaustive probe also finds a four-point network-level core and checks that
no smaller domain occurs within universes of size at most three; this bounded minimality
claim is executable evidence, not an unbounded theorem.

\section{Target D1: locally safe fallback edges can cycle globally}

\statusbox{darkred}{Result \evX: D1 is false as a local-to-global implication.
There is an exact down-spectrum-gated phase-1 run with two fallback groups, no common
model-agreeing witness in either group, and every fallback pair $\PO$, yet the declared
order contains a directed cycle.}

\subsection{A 14-point principal-ideal set model}

Let
\[
  \tau=(P\sqcup\neg P)\sqcap(Q\sqcup\neg Q),\qquad
  D=\exists\PP.\tau,\qquad C_0=D\sqcap\tau.
\]
There are no universal modalities. The formula is therefore $\forall PO$-free.

Use four minimal points $b_{00},b_{01},b_{10},b_{11}$ below both $a_1,a_2$, and the
non-transitive generating order
\[
  a_1<z_1<v_2<t_2,\qquad
  a_2<z_2<v_1<t_1,\qquad
  a_2<z_2<v_{1b}<t_{1b}.
\]
Every named point is interpreted as its principal ideal, giving the following explicit
regions and no unintended inclusions.

\begin{center}
\small
\renewcommand{\arraystretch}{1.13}
\begin{tabularx}{\linewidth}{>{\bfseries}c X >{\bfseries}c X}
\toprule
Point & Region & Point & Region \\
\midrule
$b_{00}$ & $\{0\}$ & $b_{01}$ & $\{1\}$ \\
$b_{10}$ & $\{2\}$ & $b_{11}$ & $\{3\}$ \\
$a_1$ & $\{0,1,2,3,4\}$ & $a_2$ & $\{0,1,2,3,5\}$ \\
$z_1$ & $\{0,1,2,3,4,6\}$ & $z_2$ & $\{0,1,2,3,5,7\}$ \\
$v_2$ & $\{0,1,2,3,4,6,8\}$ & $v_1$ & $\{0,1,2,3,5,7,9\}$ \\
$v_{1b}$ & $\{0,1,2,3,5,7,12\}$ & $t_2$ & $\{0,1,2,3,4,6,8,10\}$ \\
$t_1$ & $\{0,1,2,3,5,7,9,11\}$ & $t_{1b}$ & $\{0,1,2,3,5,7,12,13\}$ \\
\bottomrule
\end{tabularx}
\end{center}

The two exposed atoms assign four colors to nonmaximal points:
\[
\begin{array}{c|l}
(0,0) & b_{00},a_1,v_1,v_{1b}\\
(0,1) & b_{01},a_2,v_2\\
(1,0) & b_{10},z_1\\
(1,1) & b_{11},z_2.
\end{array}
\]
The maximal $t$ points have the displayed corresponding colors but lack $D$, since they
have no strict upper point. All nonmaximal points satisfy $D$.

Each of $a_1,a_2,v_1,v_{1b},v_2$ has all four nonmaximal closure types somewhere in its
strict lower cone. Therefore the exact refined keys satisfy
\[
  q(a_1)=q(v_1)=q(v_{1b}),\qquad q(a_2)=q(v_2).
\]

\subsection{Phase-1 generation and the cycle}

Start with the duplicate-free seed
\[
  [b_{00},b_{01},b_{10},b_{11},a_1,a_2].
\]
First-witness expansion adds $z_1,z_2$, then adds $v_2,v_1$. The latter two repeat the
exact keys of $a_2,a_1$ and are blocked; the run closes at the ten-node external set
\[
  V_0=\{b_{00},b_{01},b_{10},b_{11},a_1,a_2,z_1,z_2,v_1,v_2\}.
\]
The persistence guard $\forall\PP.D$ fails at both gate-mates, so the demand is
edge-served.

No source-model point lies above both $a_1$ and $v_1$, and none lies above both $a_2$
and $v_2$. Thus group-aware agreement is unavailable. The deterministic fallback uses
the gate children $z_1,z_2$ and declares
\[
  v_1\to z_1,\qquad v_2\to z_2.
\]
Both borrowed pairs are $\PO$ in the source model. They satisfy the full D2 pairwise
condition: the target is neither below nor disjoint from its borrower. Nevertheless,
model inclusion contains $z_1\PP v_2$ and $z_2\PP v_1$, giving
\[
  v_1\xrightarrow{B}z_1\xrightarrow{\PP}v_2
     \xrightarrow{B}z_2\xrightarrow{\PP}v_1.
\]
Hence $L^+$ is reflexive on $v_1,z_1,v_2,z_2$.

Adding $v_{1b}$ gives two borrowers sharing $z_1$ and a second cycle through the same
target. One target per group therefore does not remove cross-group cycles.

\subsection{The exact global residue}

Associate one vertex with each borrowed-target group. Draw a dependency
\[
  i\longrightarrow j
\]
when target $z_i$ is model-$\PP$-below, or equal to, some borrower in group $j$. A
declared cycle induces a cycle in this finite dependency graph; conversely, a dependency
cycle expands to a cycle alternating borrowed edges and model-$\PP$ paths. Pairwise
safety removes only self-loops. It says nothing about cycles of length at least two.

This counterexample refutes the current deterministic fallback and the inference from
per-edge safety to global irreflexivity. It does not prove that every globally coordinated
witness assignment cycles: later alternative witnesses exist in this model. A global
rank, acyclic dependency selection, feedback-edge treatment, or finite CSP remains a
possible repair.

\section{Target E: finite closure and anchor placement}

\statusbox{darkamber}{Result: E has two answers. Pure down-spectrum-gated node closure
terminates with a finite bound. The current gate-mate borrowing policy nevertheless can
produce a locally valid external fixed point that has no compatible kernel placement.
General anchor-compatible completion remains open.}

\subsection{Positive combinatorial lemma: node closure terminates}

Let $Q$ be the finite key enumeration and $\Delta$ the relevant existential
subformulas. If the gate consumes each pair $(q,\delta)\in Q\times\Delta$ at most once,
then from any finite postselected anchor seed $S_0$ it adds at most one child for each
such pair. Therefore
\[
  |E^*|\le |S_0|+|Q|\,|\Delta|
       \le |S_0|+N2^N\,|\cl(C_0)|.
\]
The same product bounds the number of productive rounds. Once $E^*$ is fixed, closing
the order and disjointness relations on that finite carrier is noniterative with respect
to point generation.

This is a set-theoretic termination statement. It is essentially the packet's reported
abstract gate theorem. It does not say that the chosen edges and kernel incidences form a
valid \texttt{MultiTier}.

\subsection{Counterexample to the current anchor borrowing policy}

Take two symmetric components $i\in\{0,1\}$ with
\[
  p_{i0}<p_{i1}<p_{i2}<v_i,\qquad p_{i2}<u_i,
  \qquad v_i\mathrel{\PO}u_i,
\]
and declare every cross-component pair $\DR$. Let $X$ hold exactly at $v_0,v_1$,
$Y$ exactly at $u_0,u_1$, and $Z$ at all six $p$ points. Put
\begin{align*}
  T &= \exists\PP.Z\sqcap\exists\PP.(X\sqcap\exists\DR.Y)
       \sqcap\exists\PP.Y,\\
  C_0 &= T\sqcap\exists\DR.T.
\end{align*}
There are no universal modalities. The first two $p$ points in each component satisfy
$C_0$ and have equal full model type. Thus $p_{i0}<p_{i1}$ is an explicit recurrence
segment, and the demand $\exists\PP.Z$ is served along the same-type $p$ chain rather
than being incidental. A component-swap automorphism gives
\[
  q(v_0)=q(v_1).
\]

The $\DR.Y$ witnesses are uniquely crossed:
\[
  v_0\xrightarrow{\DR}u_1,\qquad v_1\xrightarrow{\DR}u_0.
\]
Kernel $i$ needs both $v_i$ and $u_i$ as $\PP$ anchors above its selected phase. If
$v_0$ is the gate representative, blocked $v_1$ borrows $v_0$'s witness $u_1$, so the
external construction declares $v_1\DR u_1$. That declaration is locally harmless on
the anchor-only carrier, where there are no order edges. It cannot coexist with
$p_{11}<v_1$ and $p_{11}<u_1$, because
\[
  \DR\notin\comp(\PPI,\PP).
\]
Equivalently, downward closure of disjointness would make a point below both disjoint
from one of its strict uppers. If $v_1$ is retained first, the symmetric failure occurs
in component 0.

Thus finite key enumeration bounds how many expansion tokens are consumed, but the key
does not encode which kernel cone an anchor belongs to. The current policy can reach a
finite external frame and still fail \texttt{hkex} placement.

The example does not refute all bounded completions. An existing-witness-first policy
would choose $u_0$ for $v_1$ in this small model and avoid the clash, but no global bound
or preservation theorem is known for that policy.

\subsection{Why the old A--B family still matters}

Let
\[
  A_n=\{0,\ldots,n\},\qquad B_m=\{m+1\}.
\]
For
\[
  H=(\exists\PP.\top)\sqcap(\exists\DR.\top),\qquad
  C_0=H\sqcap\forall\PP.H,
\]
all $B_m$ have the same closure type and empty down-spectrum. Yet
\[
  \rho(A_n,B_m)=
  \begin{cases}
    \DR,&n\le m,\\
    \PPI,&n\ge m+1.
  \end{cases}
\]
For any fixed finite source-model anchor set, choose $n$ beyond every index used in it;
then every anchor is a strict part of $A_n$, and none serves $\exists\DR.\top$. Thus a
finite key space does not bound stabilization thresholds or produce a finite basis of
source-model anchors \evS.

This family does not by itself refute a fresh abstract anchor completion. Such a
completion intentionally departs from source rows and must instead prove the
anchor-compatibility lemma exposed above.

\section{Audit of the supplied probes}

\statusbox{darkred}{The supplied probes are useful exploratory artifacts, but several
headline readings are not measurements of the stated R2 construction. The original
files are preserved unchanged in the package; corrected and independent checks are kept
separately.}

\subsection{\texttt{wp131}: four material defects}

\begin{enumerate}
  \item \textbf{Wrong gate.} The target uses $q=(\mty,\text{lower spectrum})$, but
        \texttt{gated()} and gate-mate lookup in the supplied script use $\mty$ alone.
        The published 1,040-instance, 132-fallback result is therefore a type-only
        measurement, not a Target-D measurement.

  \item \textbf{The child cache omits the role.} It is keyed by $(a,\text{body})$ rather
        than $(a,\text{role},\text{body})$. If the same body occurs under two modalities,
        a $\PP$ borrower can be redirected to another role's child. In the deterministic
        supplied sweep, 18 of 1,772 borrowed targets were not $\PP$ witnesses of the
        blocker. Their relations from blocker to target were
        \[
          \DR:8,\qquad\PO:3,\qquad\EQ:6,\qquad\PPI:1.
        \]
        Adding the role changes the type-only headline counts: greedy cyclic instances
        fall from 522 to 515, and repeat-free survivors from 17 to 8.

  \item \textbf{Part G tests individuals, not groups.} The group-aware rule requires one
        witness above all borrowers in a group. Part G instead asks whether some borrower
        lacks an individual agreeing witness. The supplied predicate reports 132 fallback
        instances; the group predicate gives 135. With the actual refined key and
        relation-sensitive child cache, the corresponding counts are 164 and 167.

  \item \textbf{Part F confounds two interventions.} Its ``agree'' column calls
        \texttt{repeat\_free=True} together with \texttt{mode="agree"}. Pure group-aware
        agreement has one cycle in the supplied type-only sweep; only the combination has
        zero. The corrected refined-key finite sweep happens to have zero for both, but
        that remains randomized evidence and cannot imply D1, which the exact model in
        Section 4 refutes.
\end{enumerate}

The non-vacuity pass is also forced: sampling is conditioned on at least one borrowed
edge, then Part N accepts a class averaging at least 0.5 borrowed edges. Every retained
class averages at least one. The reported retention rates remain informative.

\subsection{\texttt{wp132}: a regression, not a placement test}

The function \texttt{fresh\_sat} returns true for every non-$\DR$ universal and false
for every existential. It therefore accepts
\[
  \forall\EQ.(P\sqcap\neg P)
\]
and rejects $\exists\EQ.P$ even when the fresh assignment makes $P$ true. It is neither
sound nor complete as a satisfiability evaluator.

The existing-point route checks only that a point is $\DR$ from a lower-cone union and
satisfies the body in the old model. It does not construct the shared top or check truth
and composition after extension. The fresh route likewise performs no RCC5 extension
check. Of the 7,525 reported refined-key group-demand evaluations, 6,038 (80.2 percent)
are accepted only by this fresh abstraction and 1,487 use an existing point. These are
correlated evaluations from 5,771 accepted generated models.

Fairly interpreted, the generated bodies are restricted to disjunctions of
$\forall\DR$ atoms, so the displayed $\EQ$ and existential errors do not directly
invalidate that narrow sweep. The probe is a purpose-built regression of local
lower-cone agreement on its generated family. It is not evidence for a realizable copied
gadget, full serviceability, or Target E.

\subsection{The wrapper masks failures}

The supplied \texttt{run\_probes.sh} executes
\begin{center}
\texttt{python3 "\$f" || echo "(non-zero exit ...)"}.
\end{center}
This converts exceptions and nonzero verdicts into shell success, defeating
\texttt{set -e}. A zero wrapper exit status is not a correctness verdict.

\section{Bounded repair directions and their exact obligations}

\subsection{No stronger existing-witness selector repairs D2}

Section 3's candidate set is a singleton. Any repair confined to choosing another
gate-mate witness is impossible. At least one of the following must change: whether $v$
is blocked, which points are available, the key/grouping discipline, or the use of the
source witness as the borrowed target.

Refusing to block $\DR$-related same-key nodes is not automatically bounded. There can
be arbitrarily many automorphic, pairwise-$\DR$ leaves with private $P$ tops. A unary
type/down-spectrum refinement does not see their multiplicity or pairwise incidence.

\subsection{A global dependency discipline for D1}

For a fixed finite external set, build the borrowed-target dependency graph from
Section 4. A globally coordinated selector can reject assignments with cyclic dependency
graphs, introduce a well-founded rank, or compute a feedback-edge set. Acyclic components
need no new D1 proof; cyclic strongly connected components are exactly where fresh
gadgets, unrolling, or kernelization are required.

This characterization is finite and checkable, but existence of an acyclic assignment
with the required labels and all frame obligations remains open.

\subsection{Fresh class-demand gadgets}

A plausible repair replaces an unsafe model witness by a fresh top
$t_{q,\delta}$ for each refined key $q$ and active $\PP$ demand $\delta$, placing every
borrower in that class below the fresh top. The packet's reported
\texttt{dkey\_union\_serves} lemma is exactly the local tool that the refuted type-only
common top lacked. If $K=N2^N$ and there are $p\le |\cl(C_0)|$ active $\PP$ demands,
the first layer uses at most $Kp$ fresh roots.

This is a bounded program, not a completed construction. Each fresh gadget has its own
demands; recursive closure can regenerate the same problem. It must also preserve
already stabilized kernel rows and package genuine recurrent suffixes as
\texttt{KernelData}. Private clones without this accounting do not prove termination.

\subsection{Conditional witness-incidence and kernel-profile keys}

If a finite witness pool $W$ is fixed in advance, refine the key to
\[
  q_W(v)=\left(q(v),\{w\in W:\rho(v,w)=\PP\}\right).
\]
Then $q_W(a)=q_W(v)$ and $\rho(a,z)=\PP$ for $z\in W$ imply
$\rho(v,z)=\PP$ literally. Every borrowed order edge agrees with the source model, so
D1 and D2 disappear. The number of keys is at most
\[
  N2^N\,2^{|W|}.
\]
The condition is circular: demand-closing $W$ can add witnesses, change the key, and
create new representatives. Target E is precisely the missing finite-pool theorem.

After kernel segments are selected, an alternative is to add each anchor's finite
kernel-incidence profile. Recording one RCC5 atom per kernel costs at most a
$5^{|\kappa|}$ factor and separates the Target-E counterexample. All abstract gate,
borrowed-edge, and frame lemmas would need to be reproved for that profile.

Weakening \texttt{ltNotDj}, simply deleting conflicting $\DR$ seeds, sending arbitrary
nonrecurrent demands to kernels, or adding unclosed private clones does not address the
actual obligations.

\section{Integrated status and recommended next steps}

The round-2 packet made real progress: the finite cross-kernel fusion theorem and the
down-spectrum cone-transfer mechanism are reported formalized. The present attack shows
that the next local selection step is not sound and that finite key enumeration does not
by itself solve semantic anchor placement.

\begin{center}
\renewcommand{\arraystretch}{1.25}
\begin{tabularx}{\linewidth}{>{\bfseries}p{0.28\linewidth}X}
\toprule
Item & Defensible status after this attack \\
\midrule
Target B / \texttt{kq\_all} & Reported Lean-certified by R2; not rebuilt from this
archive. \\
Down-spectrum local cone transfer & Reported Lean-certified and not refuted here. \\
Target D1 & Current local fallback proof route refuted by a $\PO$-only alternating
cycle. \\
Target D2 & Universal existing-witness selection refuted by a unique-witness mixed
counterexample. \\
Target E node termination & Finite combinatorial bound available. \\
Target E placement / \texttt{hkex} & Current borrowing policy refuted; general
anchor-compatible completion open. \\
Final packaging and extraction & Still open; the packet itself warns that several
``bookkeeping'' steps remain. \\
\bottomrule
\end{tabularx}
\end{center}

The most efficient research order is now:

\begin{enumerate}
  \item formalize the D1 and D2 finite countermodels as regression theorems;
  \item repair \texttt{wp131}, separate selection interventions, and make the wrapper
        propagate failures;
  \item state an anchor-compatible gate theorem with explicit kernel-incidence data;
  \item choose between a globally acyclic witness assignment, fresh class-demand
        gadgets, or a fixed-pool/profile construction, and prove its finite closure;
  \item only then rebuild \texttt{KernelData}, reindex, and assemble the final
        \texttt{MergedExtractionAt} object.
\end{enumerate}

No complete extraction theorem for the mixed quadrant, and no decidability theorem for
full $\mathcal{ALCI}_{\mathrm{RCC5}}$, is claimed.

\appendix
\begingroup
\small
\setlist[itemize]{leftmargin=1.6em,itemsep=0.1em,topsep=0.2em}
\section{Artifacts, reproducibility, and scope}

The accompanying ZIP preserves the supplied R2 packet unchanged and adds independent
checks:

\begin{itemize}
  \item \texttt{target\_d1\_multigroup\_cycle.py}: exact 14-point D1 model, phase-1
        seed expansion, refined keys, group-aware fallback, and alternating cycles;
  \item \texttt{target\_d2\_mixed\_counterexample.py}: direct mixed D2 model and
        fixed-point gate run;
  \item {\footnotesize\texttt{target\_d2\_exhaustive\_countermodel.py}}: bounded core search plus
        formula-realized and nonempty-spectrum strengthenings;
  \item \texttt{target\_e\_anchor\_compatibility.py}: symmetric kernel-anchor clash;
  \item \texttt{target\_e\_finite\_key\_anchor\_counterexample.py}: symbolic
        $A_n/B_m$ finite-anchor calculation;
  \item \texttt{wp131\_refined\_key\_audit.py}: the packet sweep with the refined key,
        relation-sensitive child cache, and group-level fallback count;
  \item \texttt{probe\_audit\_r2.py}: deterministic reproduction of the omitted
        predicates and \texttt{wp132} evaluator checks.
\end{itemize}

From the package root, run
\begin{center}
\texttt{./run\_new\_probes.sh}.
\end{center}
Every packaged file except the manifest itself is covered by
\texttt{MANIFEST.sha256}. The exact finite countermodels are independent of the
randomized seed used by the exploratory sweeps.

\subsection{Scope of the countermodels}

All exact countermodels use ordinary nonempty finite sets, so RCC5 converse and weak
composition closure follow from concrete semantics and are also checked exhaustively on
the displayed finite domains. The probes separately verify strong equality, the source
ordered-disjoint axioms, formula satisfaction, closure-type equality, exact lower
spectra, persistence guards, gate order, complete candidate sets, and the claimed final
failure.

The D1 model refutes the displayed deterministic fallback and the universal implication
from local pair safety to global acyclicity. It intentionally does not claim that every
possible simultaneous witness assignment fails. The D2 model is stronger with respect
to selection: the relevant candidate set is a singleton, so no selection among existing
gate-mate witnesses can satisfy D2. The Target-E model refutes the present borrowing
policy, while leaving open a different bounded anchor-compatible construction.

\endgroup
\end{document}
